\chapter{Esercizi: Introduzione a Python}
\label{chap:esercizi-introduzione-python}

\begin{exercise}
\label{ex:fi-python-versioni}
Che differenza c'è tra Python 2 e Python 3 e quale versione è
adottata negli appunti?
\end{exercise}

Python 3 è l'evoluzione di Python 2 e introduce modifiche non sempre
retrocompatibili, per cui un programma scritto per una versione può
richiedere adattamenti per funzionare nell'altra. Gli appunti adottano
Python 3, con esempi che richiedono la versione 3.9 o successiva
(\cref{sec:fi-python-esecuzione}).

\begin{exercise}
\label{ex:fi-python-print}
Che cosa fa la funzione \texttt{print} e perché servono le parentesi
tonde?
\end{exercise}

La funzione \texttt{print()} scrive i valori ricevuti sull'uscita standard.
Le parentesi fanno parte della sintassi di chiamata di una funzione e
racchiudono gli eventuali argomenti: per esempio,
\lstinline|print("Ciao")| stampa il messaggio \texttt{Ciao}
(\cref{sec:fi-python-esecuzione,sec:fi-python-funzioni}).

\begin{exercise}
\label{ex:fi-python-stringa}
Che cos'è una stringa?
\end{exercise}

Una \emph{stringa} è una sequenza di caratteri; in Python ha tipo
\texttt{str} e nel codice può essere delimitata da apici singoli o doppi.
Per esempio, \lstinline|"123"| è una stringa, non il numero intero
\texttt{123} (\cref{sec:fi-python-esecuzione,sec:fi-python-variabili}).

\begin{exercise}
\label{ex:fi-python-interpretato}
In che senso Python è un linguaggio interpretato?
\end{exercise}

Un programma Python viene eseguito attraverso un ambiente che ne
interpreta le operazioni, senza richiedere la produzione preventiva di
un eseguibile nativo autonomo. Nell'implementazione CPython, il sorgente
viene prima tradotto in \emph{bytecode}, poi eseguito dalla macchina
virtuale Python: compilazione e interpretazione possono quindi coesistere
(\cref{sec:fi-python-esecuzione,sec:fi-linguaggi-esecuzione}).

\begin{exercise}
\label{ex:fi-python-interattiva-edit-run}
Qual è la differenza tra modalità interattiva e ciclo
\emph{edit/run}?
\end{exercise}

Nella modalità interattiva le istruzioni vengono inserite e valutate
una alla volta nell'interprete, che mostra i risultati delle espressioni
quando previsto. Nel ciclo \emph{edit/run} si scrive o modifica un
programma in un file \texttt{.py} e lo si esegue successivamente;
le espressioni non vengono stampate automaticamente
(\cref{sec:fi-python-esecuzione}).

\begin{exercise}
\label{ex:fi-python-prompt}
Che cosa indica il prompt \lstinline|>>>| dell'interprete Python?
\end{exercise}

Il prompt \lstinline|>>>| indica che l'interprete è in modalità
interattiva ed è pronto a ricevere un'istruzione o un'espressione.
Non fa parte del codice digitato dall'utente
(\cref{sec:fi-python-esecuzione}).

\begin{exercise}
\label{ex:fi-python-variabile-matematica}
Che cos'è una variabile e in che cosa differisce da una variabile
matematica?
\end{exercise}

In Python una variabile è un nome associato a un valore durante
l'esecuzione, e l'associazione può essere modificata mediante
assegnamento. In matematica una variabile rappresenta invece un valore
in una relazione, senza che il simbolo di uguaglianza comporti un
aggiornamento dello stato (\cref{sec:fi-python-variabili}).

\begin{exercise}
\label{ex:fi-python-assegnamento-incremento}
Perché l'istruzione \texttt{x = x + 1} ha senso in programmazione?
\end{exercise}

L'operatore \texttt{=} indica un assegnamento, non un'uguaglianza
matematica. Viene prima valutata l'espressione \texttt{x + 1}
utilizzando il valore corrente di \texttt{x}; il risultato viene poi
associato allo stesso nome. L'istruzione incrementa quindi il valore
di uno, purché \texttt{x} sia già definita e l'operazione sia valida
(\cref{sec:fi-python-variabili}).

\begin{exercise}
\label{ex:fi-python-operatori-aritmetici}
Quali sono gli operatori aritmetici principali di Python e che
differenza c'è tra \texttt{/} e \lstinline|//|?
\end{exercise}

Gli operatori \texttt{+}, \texttt{-}, \texttt{*} e \texttt{**}
rappresentano rispettivamente addizione, sottrazione, moltiplicazione
e potenza. L'operatore \texttt{/} esegue la divisione ordinaria e
restituisce un \texttt{float} per operandi numerici ordinari, mentre
\lstinline|//| arrotonda il quoziente verso \(-\infty\): per esempio,
\texttt{17 / 5} vale \texttt{3.4} e \lstinline|17 // 5| vale
\texttt{3} (\cref{sec:fi-python-variabili}).

\begin{exercise}
\label{ex:fi-python-operatore-modulo}
Che cosa restituisce l'operatore \texttt{\%} e in quali situazioni può
essere utile?
\end{exercise}

L'operatore \texttt{\%} restituisce il resto della divisione:
\lstinline|17 % 5| vale \texttt{2}. È utile, per esempio, per
verificare se un intero è pari mediante \lstinline|n % 2 == 0| o
per distinguere i multipli di un intero positivo
(\cref{sec:fi-python-variabili,sec:fi-python-funzioni}).

\begin{exercise}
\label{ex:fi-python-tipi-elementari}
Quali sono i principali tipi di dato elementari introdotti
negli appunti?
\end{exercise}

I tipi fondamentali introdotti sono \texttt{int} per gli interi,
\texttt{float} per i numeri in virgola mobile, \texttt{str} per le
stringhe e \texttt{bool} per i valori logici. In Python il tipo
appartiene al valore; \texttt{type()} permette di conoscerlo
(\cref{sec:fi-python-variabili}).

\begin{exercise}
\label{ex:fi-python-stringa-numero}
Che differenza c'è tra una stringa e un numero, anche quando
contengono gli stessi caratteri?
\end{exercise}

Una stringa rappresenta testo, mentre un valore numerico può essere
impiegato direttamente nelle operazioni aritmetiche. Per esempio,
\lstinline|"3" + "4"| produce \lstinline|"34"| per concatenazione,
mentre \texttt{3 + 4} produce \texttt{7}. Per convertire una stringa
numerica in un intero si può usare \texttt{int()}
(\cref{sec:fi-python-variabili}).

\begin{exercise}
\label{ex:fi-python-input-conversione}
Che cosa fa la funzione \texttt{input} e perché spesso è necessario
convertire il suo risultato?
\end{exercise}

La funzione \texttt{input()} legge una riga dall'ingresso standard e
restituisce sempre una stringa, anche se l'utente ha digitato cifre.
Per utilizzare il dato in un calcolo numerico occorre convertirlo,
per esempio con \texttt{int()} o \texttt{float()}; una stringa non
convertibile provoca \texttt{ValueError}
(\cref{sec:fi-python-variabili}).

\begin{exercise}
\label{ex:fi-python-booleani}
Che cosa sono i valori booleani \texttt{True} e \texttt{False}?
\end{exercise}

Sono i due valori del tipo \texttt{bool}, utilizzati per esprimere
condizioni vere o false. Costituiscono, per esempio, il risultato
degli operatori di confronto e permettono di controllare il flusso
del programma (\cref{sec:fi-python-variabili,sec:fi-python-condizioni}).

\begin{exercise}
\label{ex:fi-python-operatori-logici}
Come funzionano gli operatori logici \texttt{and}, \texttt{or} e
\texttt{not}?
\end{exercise}

Per operandi booleani, \texttt{and} restituisce \texttt{True} soltanto
se entrambe le condizioni sono vere; \texttt{or} se almeno una è vera;
\texttt{not} inverte il valore logico. Gli operatori \texttt{and} e
\texttt{or} valutano il secondo operando solo quando necessario
(\emph{cortocircuito}; \cref{sec:fi-python-condizioni}).

\begin{exercise}
\label{ex:fi-python-tabella-verita}
Che cos'è una tabella di verità?
\end{exercise}

Una \emph{tabella di verità} elenca tutte le combinazioni dei valori
logici degli operandi e il risultato dell'operazione considerata.
Per due variabili booleane comprende quattro combinazioni; nella
\cref{sec:fi-python-condizioni} sono riportate le tabelle di
\texttt{and}, \texttt{or} e \texttt{not}.

\begin{exercise}
\label{ex:fi-python-sintassi-if}
Qual è la struttura sintattica di un'istruzione \texttt{if} in Python?
\end{exercise}

L'istruzione inizia con \lstinline|if condizione:|, seguita da
un blocco di istruzioni indentato. Il blocco viene eseguito soltanto
se la condizione è vera; il carattere \texttt{:} e l'indentazione
sono parte della sintassi (\cref{sec:fi-python-condizioni}).

\begin{exercise}
\label{ex:fi-python-else-elif}
A che cosa servono \texttt{else} ed \texttt{elif}?
\end{exercise}

\texttt{elif} introduce una condizione alternativa, valutata quando
quelle precedenti non sono soddisfatte; \texttt{else} indica il
blocco da eseguire se nessuna delle condizioni precedenti è vera.
In una catena \texttt{if}--\texttt{elif}--\texttt{else} viene eseguito
al più un ramo (\cref{sec:fi-python-condizioni}).

\begin{exercise}
\label{ex:fi-python-indentazione}
Perché l'indentazione è obbligatoria e significativa in Python?
\end{exercise}

L'indentazione delimita i blocchi e ne determina l'appartenenza
alle istruzioni di controllo o alle funzioni. Un rientro non valido
può causare \texttt{IndentationError}; un rientro formalmente
ammesso ma collocato nel punto sbagliato può cambiare il comportamento
del programma (\cref{sec:fi-python-condizioni}).

\begin{exercise}
\label{ex:fi-python-blocco-istruzioni}
Che cos'è un blocco di istruzioni?
\end{exercise}

È un insieme di istruzioni appartenenti alla stessa struttura,
come il corpo di una funzione, di una condizione o di un ciclo.
In Python le istruzioni dello stesso blocco condividono il livello
di indentazione; ulteriori rientri identificano eventuali blocchi
annidati (\cref{sec:fi-python-condizioni,sec:fi-python-cicli}).

\begin{exercise}
\label{ex:fi-python-cicli-condizionali}
Che differenza c'è tra cicli e istruzioni condizionali?
\end{exercise}

Un'istruzione condizionale seleziona un blocco in base a una
condizione; un ciclo ripete un blocco per più iterazioni, fino al
raggiungimento della condizione di arresto o all'esaurimento dei
valori da visitare (\cref{sec:fi-python-condizioni,sec:fi-python-cicli}).

\begin{exercise}
\label{ex:fi-python-quando-while}
Quando conviene usare un ciclo \texttt{while}?
\end{exercise}

Il ciclo \texttt{while} è appropriato quando il numero delle
iterazioni non è noto in anticipo e la ripetizione dipende dal
permanere di una condizione, per esempio durante la validazione
di un dato inserito dall'utente. La condizione viene verificata
prima di ogni iterazione (\cref{sec:fi-python-cicli}).

\begin{exercise}
\label{ex:fi-python-ciclo-infinito}
Che cos'è un ciclo infinito e perché può essere un errore grave?
\end{exercise}

È un ciclo che non raggiunge mai una condizione di arresto.
Se non è intenzionale, può impedire il completamento del programma
e consumare inutilmente risorse. Occorre verificare che gli
aggiornamenti eseguiti nel corpo possano rendere falsa la condizione
del \texttt{while} (\cref{sec:fi-python-cicli,sec:fi-algoritmi}).

\begin{exercise}
\label{ex:fi-python-for-range}
Che cosa fa il ciclo \texttt{for} e perché è utile con
\texttt{range}?
\end{exercise}

Il ciclo \texttt{for} visita nell'ordine gli elementi di un oggetto
iterabile. Con \texttt{range()} permette di iterare su una sequenza
di interi definita da estremi e, facoltativamente, da un passo,
senza aggiornare manualmente un contatore
(\cref{sec:fi-python-cicli}).

\begin{exercise}
\label{ex:fi-python-range-esempi}
Come funzionano \texttt{range(5)}, \texttt{range(1, 6)} e
\texttt{range(1, 10, 2)}?
\end{exercise}

\texttt{range(5)} genera gli interi da \(0\) a \(4\),
\texttt{range(1, 6)} quelli da \(1\) a \(5\), mentre
\texttt{range(1, 10, 2)} genera \(1,3,5,7,9\).
L'estremo iniziale è incluso, quello finale escluso e il terzo
argomento, quando presente, specifica il passo
(\cref{sec:fi-python-cicli}).

\begin{exercise}
\label{ex:fi-python-accumulatore}
Perché è utile accumulare valori dentro una variabile durante un
ciclo?
\end{exercise}

Un \emph{accumulatore} conserva il risultato parziale delle iterazioni
precedenti e viene aggiornato a ogni passaggio. Permette, per
esempio, di calcolare la somma di una sequenza senza memorizzare
tutte le somme intermedie (\cref{sec:fi-python-cicli}).

\begin{exercise}
\label{ex:fi-python-testo-binario}
Che cos'è un file di testo e in che cosa differisce da un file
binario?
\end{exercise}

In un file di testo i byte vengono interpretati come caratteri
attraverso una codifica, per esempio UTF-8. Un file aperto in
modalità binaria viene invece letto o scritto come sequenza di byte,
senza conversione automatica in caratteri
(\cref{sec:fi-python-file,sec:fi-caratteri}).

\begin{exercise}
\label{ex:fi-python-modalita-file}
Che cosa significa aprire un file in modalità \lstinline|"w"|,
\lstinline|"r"| oppure \lstinline|"a"|?
\end{exercise}

La modalità \lstinline|"w"| crea un file per la scrittura o ne
sovrascrive il contenuto, se esiste già; \lstinline|"r"| apre
un file esistente in lettura; \lstinline|"a"| permette di
aggiungere dati alla fine, creando il file se necessario.
L'apertura in \lstinline|"w"| elimina il contenuto precedente
(\cref{sec:fi-python-file}).

\begin{exercise}
\label{ex:fi-python-with-open}
Perché è consigliabile usare \lstinline|with open(...)| quando
si lavora con i file?
\end{exercise}

Il costrutto \texttt{with} garantisce la chiusura del file
all'uscita dal blocco, anche in presenza di un'eccezione.
Evita quindi di dover gestire manualmente la chiusura della
risorsa (\cref{sec:fi-python-file}).

\begin{exercise}
\label{ex:fi-python-read-for}
Che differenza c'è tra leggere un file con \texttt{read()} e
leggerlo riga per riga con un ciclo \texttt{for}?
\end{exercise}

\texttt{read()} senza argomenti restituisce l'intero contenuto
del file di testo come un'unica stringa; un ciclo \texttt{for}
lo visita invece una riga alla volta. Quest'ultima modalità è
preferibile per file grandi quando non è necessario caricare
tutto il contenuto in memoria (\cref{sec:fi-python-file}).

\begin{exercise}
\label{ex:fi-python-conversione-file}
Perché quando leggiamo numeri da un file dobbiamo usare
\texttt{int} o \texttt{float}?
\end{exercise}

La lettura di un file di testo restituisce stringhe, anche se
contengono cifre. Per eseguire operazioni numeriche è quindi
necessario convertirle in valori \texttt{int} o \texttt{float};
se il testo non è convertibile, viene sollevata un'eccezione
\texttt{ValueError} (\cref{sec:fi-python-file,sec:fi-python-variabili}).

\begin{exercise}
\label{ex:fi-python-file-inesistente}
Che cosa succede se si prova a leggere un file che non esiste e
come si può gestire il problema?
\end{exercise}

L'apertura in modalità \lstinline|"r"| di un file inesistente
solleva \texttt{FileNotFoundError}. Si può gestire l'eccezione
mediante \texttt{try} e \texttt{except FileNotFoundError},
per esempio segnalando l'assenza del file senza interrompere
bruscamente il programma (\cref{sec:fi-python-file}).
